\section{Introduction}
\label{sec:intro}

Elasticity is a core promise of cloud computing~\cite{mell2011nist}: applications should add capacity when demand rises and release it when demand falls. In many modern systems, demand does not arrive only as direct HTTP requests. It is often buffered through message queues, such as order-processing jobs during an e-commerce flash sale, telemetry events from Internet-of-Things devices, or asynchronous tasks in a machine-learning inference pipeline. In these workloads, the queue depth is an immediate measure of unfinished work. If the consumer deployment reacts late, users may experience long processing delays even though the cluster eventually scales out.

Kubernetes~\cite{kubernetes2024} provides the Horizontal Pod Autoscaler (HPA), which adjusts a Deployment's replica count so that a selected resource metric, usually CPU or memory utilisation, approaches a target value~\cite{kubernetes2024hpa}. HPA is simple and widely available, but CPU is an indirect and sometimes delayed signal for a queue-backed service. A consumer can be CPU-saturated while the backlog is either growing rapidly or almost drained; CPU alone does not tell the autoscaler how much work remains in the queue.

Kubernetes Event-Driven Autoscaling (KEDA)~\cite{keda2024} offers a different approach. Instead of asking only how busy the pods are, KEDA can ask how much work is waiting outside the pods. For RabbitMQ, KEDA can scale a consumer directly from queue length. Internally, KEDA still creates an HPA object, but the metric fed into the HPA is an external business metric rather than a CPU metric. This makes KEDA a useful comparison point for event-driven cloud workloads.

\textbf{Project nature and research question.} This is a technical cloud-computing project: we implement and evaluate existing autoscaling mechanisms rather than proposing a new autoscaling algorithm. The central question is: \emph{for the same RabbitMQ-backed consumer and the same burst workload, how do CPU-based HPA and queue-length-based KEDA differ in scale-up responsiveness, queue-drain performance, and scale-down behaviour?}

\textbf{Methodology.} We built a complete reproducibility artifact on GitHub~\cite{group5artifact2026}. The artifact deploys RabbitMQ, a Spring Boot consumer, HPA/KEDA policies, a Python load generator, and metric-collection scripts on a single-node K3s cluster. To make the comparison fair, both strategies use the same Docker image, CPU and memory requests/limits, minimum and maximum replica counts, and HPA \texttt{behavior} rules. The only independent variable is the signal used for scaling: CPU utilisation for HPA and RabbitMQ queue length for KEDA.

\textbf{Main findings.} KEDA reacts earlier, reaches useful capacity sooner, drains the queue faster, and is the only strategy that scales down during the observed post-burst window. The main reason is not that KEDA uses a more aggressive control loop. It uses a better signal for this workload. Queue length decreases continuously as work is completed, while CPU remains pinned near the pod limit as long as any backlog remains.

The rest of this report is organized as follows. Section~\ref{sec:background} reviews HPA and KEDA. Section~\ref{sec:design} presents the implemented system. Section~\ref{sec:setup} describes the experiment. Section~\ref{sec:results} analyzes the results. Section~\ref{sec:discussion} discusses implications and limitations, and Section~\ref{sec:conclusion} concludes.
